\documentclass[12pt]{article}
\usepackage[utf8]{inputenc}
\usepackage{amsmath,amssymb,amsthm}
\usepackage{algorithm}
\usepackage{algpseudocode}
\usepackage{geometry}
\geometry{margin=1in}

\newtheorem{theorem}{Theorem}
\newtheorem{lemma}[theorem]{Lemma}
\newtheorem{definition}{Definition}

\title{Problem 112: Bouncy Numbers}
\author{Project Euler Solution}
\date{}

\begin{document}
\maketitle

\section{Problem Statement}

A positive integer is \textbf{increasing} if its digits form a non-decreasing sequence, \textbf{decreasing} if non-increasing, and \textbf{bouncy} if neither. Find the least positive integer $n$ such that the proportion of bouncy numbers among $\{1,2,\ldots,n\}$ is exactly $99\%$.

\section{Mathematical Development}

\begin{definition}
Let $n$ have decimal digits $d_1 d_2 \cdots d_k$. We call $n$ \emph{increasing} if $d_i \le d_{i+1}$ for all $i$, \emph{decreasing} if $d_i \ge d_{i+1}$ for all $i$, and \emph{bouncy} if it is neither.
\end{definition}

\begin{theorem}[Small numbers are non-bouncy]
Every positive integer $n < 100$ is non-bouncy.
\end{theorem}

\begin{proof}
Single-digit numbers vacuously satisfy both monotonicity conditions. For two-digit numbers $d_1 d_2$, either $d_1 \le d_2$ or $d_1 > d_2$ by the trichotomy of the integers, so the number is increasing or decreasing.
\end{proof}

\begin{theorem}[Bouncy density]
Let $B(n)$ denote the count of bouncy integers in $\{1,\ldots,n\}$. Then $B(n)/n \to 1$ as $n \to \infty$.
\end{theorem}

\begin{proof}
The non-bouncy count below $10^k$ is $O(k^{10})$ (a polynomial in the digit count, from the stars-and-bars formula of Problem~113). The total count $10^k - 1$ is exponential in $k$. Since $k^{10}/10^k \to 0$, the non-bouncy fraction vanishes.
\end{proof}

\begin{lemma}[Divisibility constraint]
If $B(n)/n = 99/100$, then $100 \mid n$.
\end{lemma}

\begin{proof}
$100 B(n) = 99n$ implies $n = 100(n - B(n))$, and $n - B(n)$ is a positive integer, so $100 \mid n$.
\end{proof}

\begin{theorem}[Existence and computability]
There exists a least $n$ with $B(n)/n = 99/100$, and a sequential scan finds it in finite time.
\end{theorem}

\begin{proof}
By Theorem~2, $B(n)/n$ eventually exceeds $99/100$. At $n = 100$, $B(100) = 0 < 99$ (Theorem~1). Since the non-bouncy count grows polynomially while $n/100$ grows linearly, a crossing must occur. The scan examines every integer and maintains an exact count, so the first qualifying $n$ cannot be missed.
\end{proof}

\section{Algorithm}

\begin{algorithm}[H]
\caption{Find least $n$ with $B(n)/n = 99/100$}
\begin{algorithmic}[1]
\State $b \gets 0$
\For{$n = 1, 2, 3, \ldots$}
    \If{$n$ is bouncy}
        \State $b \gets b + 1$
    \EndIf
    \If{$100b = 99n$}
        \State \Return $n$
    \EndIf
\EndFor
\end{algorithmic}
\end{algorithm}

A number is bouncy if and only if its digit sequence contains both a strict increase ($d_i < d_{i+1}$ for some $i$) and a strict decrease ($d_j > d_{j+1}$ for some $j$).

\section{Complexity Analysis}

\begin{itemize}
\item \textbf{Time:} $O(N \log N)$ where $N = 1{,}587{,}000$. Each bouncy check costs $O(\log n)$ for digit extraction.
\item \textbf{Space:} $O(\log N)$ for the digit representation.
\end{itemize}

\section{Answer}

\[
\boxed{1587000}
\]

\end{document}
